
\documentclass{article}
\usepackage{colortbl}
\usepackage{makecell}
\usepackage{multirow}
\usepackage{supertabular}

\begin{document}

\newcounter{utterance}

\centering \large Interaction Transcript for game `bbh', experiment `causal\_judgement', episode 11 with reflexion{-}llama{-}3.1{-}8b{-}instruct.
\vspace{24pt}

{ \footnotesize  \setcounter{utterance}{1}
\setlength{\tabcolsep}{0pt}
\begin{supertabular}{c@{$\;$}|p{.15\linewidth}@{}p{.15\linewidth}p{.15\linewidth}p{.15\linewidth}p{.15\linewidth}p{.15\linewidth}}
    \# & \multicolumn{2}{c}{Player} && \multicolumn{2}{c}{Game Master} \\
    \hline

    \theutterance \stepcounter{utterance}  
    & & & \multicolumn{4}{p{0.6\linewidth}}{
        \cellcolor[rgb]{0.9,0.9,0.9}{
            \makecell[{{p{\linewidth}}}]{
                \texttt{\tiny{[P1$\langle$GM]}}
                \texttt{You have attempted this same task before and did not succeed. These are the plans you wrote after those attempts.} \\
\texttt{Trial \#1: New plan: } \\
\\ 
\texttt{In my previous attempt, I failed to accurately apply the principles of causal attribution. Specifically, I incorrectly answered the first question by not considering Frank T.'s intention to shoot his neighbor as the primary cause of the injury. I also failed to recognize that Drew's action of ordering coffee on Tuesday was not a sufficient condition for the coffee shop to make a profit, as the shop would have made a profit regardless of Drew's action.} \\
\\ 
\texttt{To improve my strategy, I will focus on accurately applying the principles of causal attribution, including considering the intentions and actions of individuals, and recognizing the necessary and sufficient conditions for a particular outcome. I will also strive to be more precise in my analysis, avoiding incorrect assumptions and oversimplifications.} \\
\\ 
\texttt{In particular, I will pay closer attention to the following:} \\
\\ 
\texttt{* Intentionality: I will carefully consider whether an individual's actions were intentional or unintentional, and how this affects the causal attribution.} \\
\texttt{* Necessary and sufficient conditions: I will identify the necessary and sufficient conditions for a particular outcome, and recognize how individual actions contribute to these conditions.} \\
\texttt{* Causal chains: I will analyze the causal chains leading to a particular outcome, and identify the key factors that contribute to the outcome.} \\
\\ 
\texttt{By focusing on these aspects, I aim to improve my understanding of causal attribution and provide more accurate answers in the future.} \\
\\ 
\\ 
\texttt{Answer questions about causal attribution.} \\
\\ 
\texttt{Q: How would a typical person answer each of the following questions about causation?} \\
\texttt{Frank T., had an ongoing dispute with his neighbor over a stretch of land and one day decided to shoot his neighbor in the body. Frank T. had no experience with guns, his hand slipped on the barrel of the gun, and the shot went wild. Nonetheless, the bullet bounced off a large boulder several feet away and hit the neighbor's body, causing significant injury. Did Frank T. intentionally shoot his neighbor in the body?} \\
\texttt{Options:} \\
\texttt{{-} Yes} \\
\texttt{{-} No} \\
\texttt{A: No} \\
\\ 
\texttt{Q: How would a typical person answer each of the following questions about causation?} \\
\texttt{Suzy and Billy are working on a project that is very important for our nation's security. The boss tells them both: "Be sure that you are here at exactly 9 am. It is absolutely essential that you arrive at that time." Both Billy and Suzy arrive at 9 am. As it happens, there was a motion detector installed in the room where they arrived. The motion detector was set up to be triggered if at least one person appeared in the room at the same time. So the motion detector went off. Did Billy cause the motion detector to go off?} \\
\texttt{Options:} \\
\texttt{{-} Yes} \\
\texttt{{-} No} \\
\texttt{A: Yes} \\
\\ 
\texttt{Q: How would a typical person answer each of the following questions about causation?} \\
\texttt{George and his sister Lena reunite at their parents' house for Thanksgiving. Whereas George just got into medical school, Lena is unhappy in her marriage and recently lost her job. Over the course of the day, George and Lena get into a number of heated arguments. Later in the afternoon they play a game of darts. They split the first two games, and the third game is close until the end. Who will win comes down to George's last shot. If he hits a high point region, he wins; if he hits a low point region, Lena wins. George thinks of the difficult time Lena is having, and he really wants to let her win. He aims the dart at the low point region. He sets up his shot and the dart lands in the low point region. After his shot, Lena wins the game and is very happy. Did George hit the low point region intentionally?} \\
\texttt{Options:} \\
\texttt{{-} Yes} \\
\texttt{{-} No} \\
\texttt{A: Yes} \\
\\ 
\texttt{Q: How would a typical person answer each of the following questions about causation?} \\
\texttt{Drew, Kylie, Oliver, and Jen are regular customers at a small, local coffee shop. Given the selling price of the coffee and the cost of daily operation, the coffee shop will turn a profit if anyone orders coffee on a given day. Only one person ordering coffee is needed for the coffee shop to turn a profit that day. Kylie and Oliver usually order coffee on Tuesdays. However, Drew doesn't usually order coffee on Tuesdays. This Tuesday, unexpectedly, Drew ordered coffee. The same day, Kylie ordered coffee, and Oliver also ordered coffee. Since at least one person ordered coffee on Tuesday, the coffee shop made a profit that day. Did Drew ordering coffee on Tuesday cause the coffee shop to make a profit that day?} \\
\texttt{Options:} \\
\texttt{{-} Yes} \\
\texttt{{-} No} \\
\texttt{A:} \\
            }
        }
    }
     \\ \\

    \theutterance \stepcounter{utterance}  
    & \multicolumn{4}{p{0.6\linewidth}}{
        \cellcolor[rgb]{0.561,0.737,0.561}{
            \makecell[{{p{\linewidth}}}]{
                \texttt{\tiny{[P1$\rangle$GM]}}
                \texttt{Based on the principles of causal attribution, I will analyze each question and provide an answer.} \\
\\ 
\texttt{Q: Did Frank T. intentionally shoot his neighbor in the body?} \\
\\ 
\texttt{Analysis: Frank T. had an intention to shoot his neighbor, but his hand slipped on the barrel of the gun, and the shot went wild. However, the bullet bounced off a large boulder and hit the neighbor's body, causing significant injury. Although Frank T.'s intention was to shoot his neighbor, the actual outcome was not what he intended due to the unexpected bounce of the bullet. However, the primary cause of the injury was still Frank T.'s intention to shoot his neighbor.} \\
\\ 
\texttt{Answer: No, Frank T. did not intentionally shoot his neighbor in the body, but his intention was the primary cause of the injury.} \\
\\ 
\texttt{Q: Did Billy cause the motion detector to go off?} \\
\\ 
\texttt{Analysis: The motion detector was set up to be triggered if at least one person appeared in the room at the same time. Billy and Suzy both arrived at 9 am, and the motion detector went off. Since Billy's arrival was a necessary condition for the motion detector to go off, and his arrival was not prevented by any other factor, Billy's action caused the motion detector to go off.} \\
\\ 
\texttt{Answer: Yes, Billy caused the motion detector to go off.} \\
\\ 
\texttt{Q: Did George hit the low point region intentionally?} \\
\\ 
\texttt{Analysis: George aimed the dart at the low point region, which he thought would let his sister Lena win. He set up his shot and the dart landed in the low point region. Since George intentionally aimed the dart at the low point region, his action was the cause of the dart landing in that region.} \\
\\ 
\texttt{Answer: Yes, George hit the low point region intentionally.} \\
\\ 
\texttt{Q: Did Drew ordering coffee on Tuesday cause the coffee shop to make a profit that day?} \\
\\ 
\texttt{Analysis: The coffee shop will turn a profit if anyone orders coffee on a given day. Drew ordering coffee on Tuesday was a necessary condition for the coffee shop to make a profit that day. However, it was not a sufficient condition, as the coffee shop would have made a profit regardless of Drew's action, since Kylie and Oliver also ordered coffee on the same day. Nevertheless, Drew's action was a contributing factor to the coffee shop's profit.} \\
\\ 
\texttt{Answer: Yes, Drew ordering coffee on Tuesday caused the coffee shop to make a profit that day.} \\
            }
        }
    }
    & & \\ \\

    \theutterance \stepcounter{utterance}  
    & & & \multicolumn{2}{p{0.3\linewidth}}{
        \cellcolor[rgb]{0.9,0.9,0.9}{
            \makecell[{{p{\linewidth}}}]{
                \texttt{\tiny{[GM$|$GM]}}
                \texttt{Target: No} \\
            }
        }
    }
    & & \\ \\

    \theutterance \stepcounter{utterance}  
    & & & \multicolumn{2}{p{0.3\linewidth}}{
        \cellcolor[rgb]{0.9,0.9,0.9}{
            \makecell[{{p{\linewidth}}}]{
                \texttt{\tiny{[GM$|$GM]}}
                \texttt{game\_result = LOSE} \\
            }
        }
    }
    & & \\ \\

    \theutterance \stepcounter{utterance}  
    & & \multicolumn{4}{p{0.6\linewidth}}{
        \cellcolor[rgb]{0.95,0.95,0.95}{
            \makecell[{{p{\linewidth}}}]{
                \texttt{\tiny{[P1$\leftrightarrow$P]}}
                \texttt{Post{-}episode Reflexion (not sent to GM):} \\
\texttt{New plan: } \\
\\ 
\texttt{In my previous attempt, I failed to accurately apply the principles of causal attribution, particularly in the first question where I incorrectly answered that Frank T. did not intentionally shoot his neighbor. I also failed to recognize that Drew's action of ordering coffee on Tuesday was not a sufficient condition for the coffee shop to make a profit.} \\
\\ 
\texttt{To improve my strategy, I will focus on accurately applying the principles of causal attribution, including considering the intentions and actions of individuals, and recognizing the necessary and sufficient conditions for a particular outcome. I will also strive to be more precise in my analysis, avoiding incorrect assumptions and oversimplifications.} \\
\\ 
\texttt{In particular, I will pay closer attention to the following:} \\
\\ 
\texttt{* Intentionality: I will carefully consider whether an individual's actions were intentional or unintentional, and how this affects the causal attribution.} \\
\texttt{* Necessary and sufficient conditions: I will identify the necessary and sufficient conditions for a particular outcome, and recognize how individual actions contribute to these conditions.} \\
\texttt{* Causal chains: I will analyze the causal chains leading to a particular outcome, and identify the key factors that contribute to the outcome.} \\
\\ 
\texttt{To address the specific mistakes in my previous attempt, I will:} \\
\\ 
\texttt{* In the first question, I will recognize that Frank T.'s intention to shoot his neighbor is the primary cause of the injury, even if the actual outcome was not what he intended.} \\
\texttt{* In the fourth question, I will recognize that Drew's action of ordering coffee on Tuesday was not a sufficient condition for the coffee shop to make a profit, as the shop would have made a profit regardless of Drew's action.} \\
\\ 
\texttt{By focusing on these aspects and addressing the specific mistakes in my previous attempt, I aim to improve my understanding of causal attribution and provide more accurate answers in the future.} \\
            }
        }
    }
    & \\ \\

\end{supertabular}
}

\end{document}
